\documentclass[12pt, a4paper]{article}

% 引入必要的宏包
\usepackage[UTF8]{ctex}     % 中文支持
\usepackage{amsmath, amssymb, amsthm} % 数学符号和环境
\usepackage{geometry}       % 页面设置
\usepackage{enumitem}       % 列表定制
\usepackage{fancyhdr}       % 页眉页脚
\usepackage{cases}          % 分段函数增强

% 页面边距设置
\geometry{left=2.5cm, right=2.5cm, top=2.5cm, bottom=2.5cm}

% 宏定义：方便排版选择题选项
% 横向排列的选项 (4列)
\newcommand{\choicesFour}[4]{
    \begin{enumerate}[label=\Alph*., itemsep=0pt, parsep=0pt, topsep=5pt, labelsep=0.5em]
        \begin{minipage}{0.22\linewidth} \item #1 \end{minipage}
        \begin{minipage}{0.22\linewidth} \item #2 \end{minipage}
        \begin{minipage}{0.22\linewidth} \item #3 \end{minipage}
        \begin{minipage}{0.22\linewidth} \item #4 \end{minipage}
    \end{enumerate}
}
% 横向排列的选项 (2列)
\newcommand{\choicesTwo}[4]{
    \begin{enumerate}[label=\Alph*., itemsep=0pt, parsep=0pt, topsep=5pt, labelsep=0.5em]
        \begin{minipage}{0.45\linewidth} \item #1 \end{minipage}
        \begin{minipage}{0.45\linewidth} \item #2 \end{minipage}
        \begin{minipage}{0.45\linewidth} \item #3 \end{minipage}
        \begin{minipage}{0.45\linewidth} \item #4 \end{minipage}
    \end{enumerate}
}
% 纵向排列的选项 (1列)
\newcommand{\choices}[4]{
    \begin{enumerate}[label=\Alph*., itemsep=0pt, parsep=0pt, topsep=5pt, labelsep=0.5em]
        \item #1
        \item #2
        \item #3
        \item #4
    \end{enumerate}
}

% 填空题下划线
\newcommand{\blank}[1]{\underline{\hspace{#1}}}

\title{\textbf{《高等数学》必刷习题——提高版}}
\author{}
\date{}

\begin{document}

\section*{第六章 不定积分（提高）}

\begin{enumerate}
    \item 设函数 $ f(x) $ 的一个原函数为 $ \frac{1}{x} $ ，则 $ f(x)= $ \blank{1cm}。
    \choicesFour{$ -\frac{1}{x^{2}} $}{$ \frac{2}{x^{3}} $}{$ \frac{1}{x} $}{$ \ln|x| $}

    \item 若 $ e^{-x} $ 是 $ f(x) $ 的一个原函数，则 $ \int xf(x)dx = $ \blank{1cm}。
    \choicesTwo
    {$ e^{-x}(1-x)+C $}
    {$ e^{-x}(x+1)+C $}
    {$ e^{-x}(x-1)+C $}
    {$ -e^{-x}(x+1)+C $}

    \item 设 $ \int f(x)dx = \ln(x + \sqrt{1 + x^{2}}) + C, $ 则 $ f(x) = $ \blank{1cm}。
    \choicesFour{$ \frac{x}{\sqrt{1+x^{2}}} $}{$ \frac{1}{\sqrt{1+x^{2}}} $}{$ \frac{1}{1+x} $}{$ \frac{1}{1+x^{2}} $}

    \item $ \int 2^{3x} dx = $ \blank{1cm}。
    \choicesFour{$ \frac{2^{3x}}{3\ln 2} + C $}{$ \frac{1}{3}(\ln 2)2^{3x} + C $}{$ \frac{2^{x}}{3} + C $}{$ \frac{2^{3x}}{\ln 2} + C $}

    \item 设函数 $ f(x) $ 在 $ [a, b] $ 上连续， $ F(x) = \int_{0}^{x^{3}} f(t) dt $ ，则 $ F'(x) = $ \blank{1cm}。
    \choicesFour{$ f(x) $}{$ f(x^{3}) $}{$ 3x^{2}f(x^{3}) $}{$ 3x^{2}f(x) $}

    \item 设 $ f'(x^{2}) = \frac{1}{x} (x > 0) $ ，则 $ f(x) = $ \blank{2cm}。

    \item 设 $ f(x) $ 的一个原函数是 $ \sin x, $ 则 $ \int x f'(x)dx = $ \blank{2cm}。

    \item $ \int 3^{x}e^{x}dx = $ \blank{2cm}。

    \item 不定积分 $ \int\frac{x-3}{x^{2}+1}dx= $ \blank{2cm}。

    \item 不定积分 $ \int \frac{1}{e^{x} + e^{-x}} dx = $ \blank{2cm}。

    \item 已知 $ f(x) $ 的一个原函数为 $ \frac{\sin x}{1 + x \sin x} $ ，求 $ \int f(x) f'(x) dx $。

    \item 若 $ \int xf(x)dx = \arcsin x + C $ 求 $ I = \int \frac{1}{f(x)}dx $。

    \item 求不定积分 $ \int\frac{e^{\sqrt{x}}\sin\sqrt{x}}{2\sqrt{x}}dx $。

    \item 计算 $ \int \sin^{2} x \cos^{3} x dx $。

    \item 计算 $ \int \frac{\sqrt{\arcsin x}}{\sqrt{1-x^{2}}} dx $。
\end{enumerate}

\end{document}
